% Options for packages loaded elsewhere
\PassOptionsToPackage{unicode}{hyperref}
\PassOptionsToPackage{hyphens}{url}
\PassOptionsToPackage{dvipsnames,svgnames,x11names}{xcolor}
\PassOptionsToPackage{space}{xeCJK}
\documentclass[
]{article}
\usepackage{xcolor}
\usepackage[margin=2.5cm]{geometry}
\usepackage{amsmath,amssymb}
\setcounter{secnumdepth}{-\maxdimen} % remove section numbering
\usepackage{iftex}
\ifPDFTeX
  \usepackage[T1]{fontenc}
  \usepackage[utf8]{inputenc}
  \usepackage{textcomp} % provide euro and other symbols
\else % if luatex or xetex
  \usepackage{unicode-math} % this also loads fontspec
  \defaultfontfeatures{Scale=MatchLowercase}
  \defaultfontfeatures[\rmfamily]{Ligatures=TeX,Scale=1}
\fi
\usepackage{lmodern}
\ifPDFTeX\else
  % xetex/luatex font selection
  \setmainfont[]{DejaVu Sans}
  \setmonofont[]{DejaVu Sans Mono}
  \ifXeTeX
    \usepackage{xeCJK}
    \setCJKmainfont[]{Noto Sans CJK SC}
  \fi
  \ifLuaTeX
    \usepackage[]{luatexja-fontspec}
    \setmainjfont[]{Noto Sans CJK SC}
  \fi
\fi
% Use upquote if available, for straight quotes in verbatim environments
\IfFileExists{upquote.sty}{\usepackage{upquote}}{}
\IfFileExists{microtype.sty}{% use microtype if available
  \usepackage[]{microtype}
  \UseMicrotypeSet[protrusion]{basicmath} % disable protrusion for tt fonts
}{}
\makeatletter
\@ifundefined{KOMAClassName}{% if non-KOMA class
  \IfFileExists{parskip.sty}{%
    \usepackage{parskip}
  }{% else
    \setlength{\parindent}{0pt}
    \setlength{\parskip}{6pt plus 2pt minus 1pt}}
}{% if KOMA class
  \KOMAoptions{parskip=half}}
\makeatother
\usepackage{longtable,booktabs,array}
\usepackage{caption}
\captionsetup[table]{skip=6pt}
\newcounter{none} % for unnumbered tables
\usepackage{calc} % for calculating minipage widths
% Correct order of tables after \paragraph or \subparagraph
\usepackage{etoolbox}
\makeatletter
\patchcmd\longtable{\par}{\if@noskipsec\mbox{}\fi\par}{}{}
\makeatother
% Allow footnotes in longtable head/foot
\IfFileExists{footnotehyper.sty}{\usepackage{footnotehyper}}{\usepackage{footnote}}
\makesavenoteenv{longtable}
\setlength{\emergencystretch}{3em} % prevent overfull lines
\providecommand{\tightlist}{%
  \setlength{\itemsep}{0pt}\setlength{\parskip}{0pt}}
\usepackage{bookmark}
\IfFileExists{xurl.sty}{\usepackage{xurl}}{} % add URL line breaks if available
\urlstyle{same}
% fallback for those not using the hyperref driver hyperxmp:
\makeatletter
\@ifundefined{xmpquote}{\newcommand{\xmpquote}[1]{#1}}{}
\makeatother
\hypersetup{
  colorlinks=true,
  linkcolor={Maroon},
  filecolor={Maroon},
  citecolor={Blue},
  urlcolor={Blue},
  pdfcreator={LaTeX via pandoc}}

\author{}
\date{}

\begin{document}

\section{筑基篇课后习题 I：P vs NP 在 pat
视角下的详细展开和论述}\label{ux7b51ux57faux7bc7ux8bfeux540eux4e60ux9898-ip-vs-np-ux5728-pat-ux89c6ux89d2ux4e0bux7684ux8be6ux7ec6ux5c55ux5f00ux548cux8bbaux8ff0}

\begin{quote}
\textbf{声明 (Declaration)}: 本文\textbf{不代表任何数学结论}
(non-mathematical-claim), 仅为基于 Lean 4
形式化的一种\textbf{实验观测报告} (experimental observation report)。

{\def\LTcaptype{none} % do not increment counter
\begin{longtable}[]{@{}
  >{\raggedright\arraybackslash}p{(\linewidth - 6\tabcolsep) * \real{0.2500}}
  >{\raggedright\arraybackslash}p{(\linewidth - 6\tabcolsep) * \real{0.2500}}
  >{\raggedright\arraybackslash}p{(\linewidth - 6\tabcolsep) * \real{0.2500}}
  >{\raggedright\arraybackslash}p{(\linewidth - 6\tabcolsep) * \real{0.2500}}@{}}
\toprule\noalign{}
\begin{minipage}[b]{\linewidth}\raggedright
习题
\end{minipage} & \begin{minipage}[b]{\linewidth}\raggedright
主题
\end{minipage} & \begin{minipage}[b]{\linewidth}\raggedright
Lean 形式化
\end{minipage} & \begin{minipage}[b]{\linewidth}\raggedright
DOI
\end{minipage} \\
\midrule\noalign{}
\endhead
\bottomrule\noalign{}
\endlastfoot
exercise\_i & P vs NP 在 pat 视角（R157） & PatPvsNPExercise.lean &
10.5281/zenodo.21916831 \\
\end{longtable}
}
\end{quote}

\textbf{Exercise I for the Foundation-Building Chapter: P vs NP Expanded
and Discussed from the PAT Perspective}

\begin{quote}
习题定位：不是算器神魂论的独立篇章，是\textbf{筑基篇的课后习题}------用筑基篇已证定理（R136--R154）把
R152/R153 的 P vs NP
侧写展开为详细论述。本题只有一个论点，不枚举。全部新定理只锚筑基篇定理，不用外部引理（延续
R153 用户纠正精神）。
\end{quote}

\emph{2026-08-13 · Internal research exercise · Lean 4 / mathlib v4.32.2
· 4 new theorems (R157), all PROVED, no sorry · 算器神魂论筑基篇课后习题
I }

\begin{center}\rule{0.5\linewidth}{0.5pt}\end{center}

\subsection{习题陈述}\label{ux4e60ux9898ux9648ux8ff0}

\begin{quote}
用筑基篇已证定理（R136--R154），把 P vs NP（R152 框架侧写 + R153 N =
相位锁定外推）展开为 pat 视角下的详细论述，并在 Lean
中形式化验证。要求：全部新定理只锚筑基篇定理，不用外部引理。
\end{quote}

\textbf{本题唯一论点：NP 的 pat 语义 = 用模糊换速度，用结构找回精确。}

\begin{itemize}
\tightlist
\item
  \textbf{模糊} = N 的存在性：witness
  的猜测不穷举------存在性由相位锁定外推给出（R150
  王氏定理），不搜遍候选。
\item
  \textbf{结构} = 验证表：有限域函数 = 预计算表（RulerLookup/R057
  一切皆表），表是结构；witness 存在 ⟺
  表条目存在（双向）------模糊的猜测经表判等还原为精确答案。
\item
  \textbf{有限域上 P = NP
  平凡}（R152）：结构已完备（表全量），模糊不必要------求解与验证都
  O(1)。
\end{itemize}

\begin{center}\rule{0.5\linewidth}{0.5pt}\end{center}

\subsection{零、总论：为什么本题只有一个论点}\label{ux96f6ux603bux8bbaux4e3aux4ec0ux4e48ux672cux9898ux53eaux6709ux4e00ux4e2aux8bbaux70b9}

R152/R153 给了 P vs NP 的侧写------\textbf{有限域上 P = NP
平凡（一切皆表）+ 验证免费 + N =
相位锁定外推}。展开它不需要枚举若干''视角''：筑基篇的动力学只有一个动作------\textbf{发散/收敛互锁对}（R047/R147）。P
vs NP
是这个动作在计算问题上的投影（复杂性理论的经典结构与之对应：``验证便宜/求解贵''的直觉
= 互锁对的两端，Aaronson {[}C12{]}、Fortnow {[}C10{]}；``计算 =
数学结构'' = 一切皆表的表结构，Wigderson {[}C9{]}）：

\begin{itemize}
\tightlist
\item
  \textbf{N（非确定性）}是发散方向的自由：猜测什么都可以，存在性不依赖搜索路径------用模糊换速度。
\item
  \textbf{验证}是收敛方向的锁定：把模糊的猜测经表判等还原为确定的真值------用结构找回精确。
\item
  \textbf{P（确定性）}是方向锁定后的唯一链（R153
  ①/R050）------没有模糊，也不需要找回。
\end{itemize}

四层论述都是同一个互锁对的展开。以下每节只补一个承载论证的定理，不堆砌。引用标注
{[}C\#{]}/{[}M\#{]}/{[}L\#{]}/{[}Q\#{]}
对应\texttt{../shared/pat\_pnp\_shared\_references\_ZH.md}（共享引用清单）
全字段清单。

\begin{center}\rule{0.5\linewidth}{0.5pt}\end{center}

\subsection{一、用结构找回精确：验证 =
查表判等（RulerLookup/RulerRevLookup）}\label{ux4e00ux7528ux7ed3ux6784ux627eux56deux7cbeux786eux9a8cux8bc1-ux67e5ux8868ux5224ux7b49rulerlookuprulerrevlookup}

\textbf{★核心：验证不重算，只判等------y = f x ⟺ (x, y) ∈ 表，双向。}

\subsubsection{论证}\label{ux8bbaux8bc1}

\begin{enumerate}
\def\labelenumi{\arabic{enumi}.}
\tightlist
\item
  \textbf{表 = 结构}（RulerLookup function\_is\_table /
  R057）：有限域上任意函数 f = 预计算表 \{(x, f
  x)\}。表是''把计算凝固成位置''的结构。对照：计算 =
  数学结构的纲领（Wigderson {[}C9{]}）、算法 = 表/结构的经典实践（Knuth
  {[}C20{]}）。
\item
  \textbf{验证 = 查表判等}（R152 ② 展开）：验证 y 是否是 f x 的答案 =
  检查 (x, y) 是否在表中------双向等价：y = f x ⟹
  表条目（lookup\_exists：表全量），表条目 ⟹ y = f
  x（lookup\_correct：表值 =
  函数值）。对照：验证便宜/求解贵的经典不对称直觉（Aaronson
  {[}C12{]}、Fortnow {[}C10{]}）------在本框架中，验证 = 查表判等，求解
  = 同一条目的前向查询（R152 有限域平凡）。
\item
  \textbf{这是''找回精确''}：模糊的候选 y
  经''表中判等''这一步还原为精确的真假------结构（表）是精确的锚。对照：交互证明中验证者的角色（Goldwasser--Micali--Rackoff
  {[}C15{]}）；信息 = 结构、可逆计算 = 信息保持（Shannon
  {[}M10{]}、Landauer {[}M11{]}、Bennett {[}M12{]}）。
\end{enumerate}

\subsubsection{形式化（PatPvsNPExercise.lean，1
定理）}\label{ux5f62ux5f0fux5316patpvsnpexercise.lean1-ux5b9aux7406}

\begin{itemize}
\tightlist
\item
  \texttt{lookup\_iff\_value}：y = f x ⟺ (x, y) ∈ makeTable
  f------双向（验证 = 查表判等）。
\end{itemize}

\textbf{验证}：一次 build 通过，0 sorry。类型需要 {[}DecidableEq
E{]}（值域可判等，验证的前提）。

\subsubsection{连接}\label{ux8fdeux63a5}

R152 的''验证 =
可逆查表后向免费''在此展开为双向判等。\textbf{验证是结构操作，不是重算}------这正是''用结构找回精确''的机制。

\begin{center}\rule{0.5\linewidth}{0.5pt}\end{center}

\subsection{二、模糊与结构等价：witness ⟺ 表条目（双向，强化
R153）}\label{ux4e8cux6a21ux7ccaux4e0eux7ed3ux6784ux7b49ux4ef7witness-ux8868ux6761ux76eeux53ccux5411ux5f3aux5316-r153}

\textbf{★核心：模糊的 witness 猜测与结构的表条目完全等价------双向。}

\subsubsection{论证}\label{ux8bbaux8bc1-1}

\begin{enumerate}
\def\labelenumi{\arabic{enumi}.}
\tightlist
\item
  \textbf{R153 只给了单向}：witness 存在 ⟹
  表条目存在（np\_witness\_in\_table）。
\item
  \textbf{反向由表结构封死}：表条目 (w, true) ∈ 表 ⟹ true = V x
  w（lookup\_correct：表值 = 函数值）⟹ V x w = true。
\item
  \textbf{双向 = 模糊/结构等价}：∃ witness（模糊）⟺ ∃
  表条目（结构）------猜测的''存在''和结构的''存在''是同一个存在。模糊没有损失，因为它和结构一一对应。对照：NP
  存在性 = 验证表条目的标准语义（Cook {[}C1{]}、Levin
  {[}C2{]}、Goldreich {[}C8{]}）------经典复杂性理论以''存在 witness +
  多项式验证''定义 NP，本框架把验证精确化为查表判等。
\end{enumerate}

\subsubsection{形式化（PatPvsNPExercise.lean，1
定理）}\label{ux5f62ux5f0fux5316patpvsnpexercise.lean1-ux5b9aux7406-1}

\begin{itemize}
\tightlist
\item
  \texttt{witness\_iff\_table\_entry}：∃ w, V x w = true ⟺ ∃ w, (w,
  true) ∈ makeTable (fun w =\textgreater{} V x w)------双向。
\end{itemize}

\textbf{验证}：0 sorry。反向的关键：\texttt{lookup\_correct} 说表值 =
函数值。教训：\texttt{fun\ w\ =\textgreater{}\ V\ x\ w} 内层 binder
与外层 w 同名时 \texttt{rw} 失效，用 \texttt{simpa} 解。

\subsubsection{连接}\label{ux8fdeux63a5-1}

这一节是 R153
的完成：\textbf{非确定性的存在性（模糊）与验证表的条目（结构）不是两个事实，是一个事实的两个投影}。

\begin{center}\rule{0.5\linewidth}{0.5pt}\end{center}

\subsection{三、★NP
的模糊/结构对偶（本题核心定理）}\label{ux4e09np-ux7684ux6a21ux7ccaux7ed3ux6784ux5bf9ux5076ux672cux9898ux6838ux5fc3ux5b9aux7406}

\textbf{★核心：NP 的 pat 语义 = 模糊存在性 ∧ 结构判等同时成立。}

\subsubsection{论证}\label{ux8bbaux8bc1-2}

\begin{enumerate}
\def\labelenumi{\arabic{enumi}.}
\tightlist
\item
  \textbf{模糊换速度}：∃ w, V x w = true------witness
  的存在性不穷举，由相位锁定外推保证（R150，见第四节）。对照：平均情形/随机化''用概率换效率''的经典路线（Rabin
  {[}C14{]}、Impagliazzo
  {[}C13{]}）------本框架的模糊是相位层的，非概率层的。
\item
  \textbf{结构找回精确}：∀ w, V x w = true ⟺ (w, true) ∈
  表------任何候选都能经表判等得到精确真值。对照：去随机化（用结构消除概率，Vadhan
  {[}C19{]}）；可逆计算（用结构保持信息，Bennett
  {[}M12{]}）------都是''结构找回精确''的经典实例。
\item
  \textbf{对偶定理}：模糊存在性 ⟹ 结构存在性 ∧
  逐点判等------猜测与判等同一互锁对（R147：求解/验证 =
  发散/收敛互锁对），发散端给出存在性，收敛端给出精确性。对照：几何相位
  = 相位锁定（Berry {[}Q1{]}、Aharonov--Bohm
  {[}Q4{]}）------相位差是物理实在，不是计算伪影。
\end{enumerate}

\subsubsection{形式化（PatPvsNPExercise.lean，1
定理）}\label{ux5f62ux5f0fux5316patpvsnpexercise.lean1-ux5b9aux7406-2}

\begin{itemize}
\tightlist
\item
  \texttt{np\_fuzzy\_structure\_duality}：∃ w, V x w = true ⟹ (∃ 表条目)
  ∧ (∀ w, V x w = true ⟺ 表条目)------★NP 的模糊/结构对偶。
\end{itemize}

\textbf{验证}：0 sorry。前半 = witness\_iff\_table\_entry 正向，后半 =
lookup\_iff\_value 逐点。

\subsubsection{连接}\label{ux8fdeux63a5-2}

本题唯一论点的定理化：\textbf{NP =
用模糊换速度（存在性），用结构找回精确（判等）}。R153 的''NP 存在性 =
验证表条目''在此从单向变成对偶。

\begin{center}\rule{0.5\linewidth}{0.5pt}\end{center}

\subsection{四、全景：模糊（N）+ 结构（验证）+
锁定（P）}\label{ux56dbux5168ux666fux6a21ux7ccan-ux7ed3ux6784ux9a8cux8bc1-ux9501ux5b9ap}

\textbf{★核心：P vs NP 在 pat 视角 = 三个结构位------N
是相位锁定外推（模糊），验证是查表判等（结构），P
是锁定方向唯一链（无模糊）。}

\subsubsection{论证}\label{ux8bbaux8bc1-3}

\begin{enumerate}
\def\labelenumi{\arabic{enumi}.}
\tightlist
\item
  \textbf{N = 相位锁定外推}（R150 王氏定理 / R153 ③）：任意未锁定相位 θ
  被 pat 格点任意精度统一（∀ ε, ∃ 格点 \textbar θ - x\textbar{} ≤
  ε）------非确定性的存在性由相位锁定外推给出，不穷举。\textbf{这是模糊换速度的存在性来源}。对照：凝聚数学（Scholze
  {[}M21{]}）以相似方式用稠密格点统一连续结构。
\item
  \textbf{验证 = 查表判等}（第一节）：∀ w, V x w = true ⟺
  表条目------结构找回精确。对照：Lean 形式化 =
  当代''用结构找回精确''的工程（de Moura {[}M28{]}、mathlib
  {[}M29{]}、Buzzard {[}M30{]}、Avigad {[}M31{]}、Hales
  {[}M23{]}）------本习题的 Lean 验收即同一工程的一个实例。
\item
  \textbf{P = 锁定方向唯一链}（R153 ① / R050 机制）：x ↦ x + d
  单射------确定性没有模糊。
\item
  \textbf{有限域上 P = NP 平凡}（R152）：表全量 ⟹ 结构已完备 ⟹
  模糊不必要------求解与验证都 O(1)，多项式与常数无别。对照：可计算性 =
  形式结构边界的经典结果（Gödel {[}M1{]}、Turing {[}M2{]}、Church
  {[}M3{]}）------有限域上一切皆表 = 结构完备即无模糊。
\item
  \textbf{诚实边界}：以上全部为结构侧写，非 P≠NP
  判定（规模增长下界不在框架内，R152 边界延续；经典复杂性层级见
  Hartmanis--Stearns {[}C4{]}、Sipser {[}C5{]}、Arora--Barak
  {[}C7{]}）。
\end{enumerate}

\subsubsection{形式化（PatPvsNPExercise.lean，1
定理）}\label{ux5f62ux5f0fux5316patpvsnpexercise.lean1-ux5b9aux7406-3}

\begin{itemize}
\tightlist
\item
  \texttt{p\_np\_pat\_perspective}：全景------(N = 外推, ∀ε ∃格点) ∧
  (验证 = 判等, ∀w) ∧ (P = 锁定链, ∀d Injective)。
\end{itemize}

\textbf{验证}：0 sorry。三个合取项分别锚 R150 / lookup\_iff\_value /
R153 ①。

\subsubsection{连接}\label{ux8fdeux63a5-3}

本题回答：\textbf{pat 视角下 P vs NP = 锁定链 vs 未锁定相位}------P
是方向锁定的唯一链（无模糊），N
是未锁定相位的存在性（模糊，王氏定理外推），验证是查表判等（结构找回精确）。这是结构论述，不是
P≠NP 判定。

\begin{center}\rule{0.5\linewidth}{0.5pt}\end{center}

\subsection{五、习题解答总结}\label{ux4e94ux4e60ux9898ux89e3ux7b54ux603bux7ed3}

{\def\LTcaptype{none} % do not increment counter
\begin{longtable}[]{@{}
  >{\raggedright\arraybackslash}p{(\linewidth - 6\tabcolsep) * \real{0.2500}}
  >{\raggedright\arraybackslash}p{(\linewidth - 6\tabcolsep) * \real{0.2500}}
  >{\raggedright\arraybackslash}p{(\linewidth - 6\tabcolsep) * \real{0.2500}}
  >{\raggedright\arraybackslash}p{(\linewidth - 6\tabcolsep) * \real{0.2500}}@{}}
\toprule\noalign{}
\begin{minipage}[b]{\linewidth}\raggedright
论证位置
\end{minipage} & \begin{minipage}[b]{\linewidth}\raggedright
内容
\end{minipage} & \begin{minipage}[b]{\linewidth}\raggedright
锚到
\end{minipage} & \begin{minipage}[b]{\linewidth}\raggedright
定理
\end{minipage} \\
\midrule\noalign{}
\endhead
\bottomrule\noalign{}
\endlastfoot
结构 & 验证 = 查表判等（双向） & RulerLookup/RulerRevLookup/R152 &
lookup\_iff\_value \\
模糊↔结构 & witness ⟺ 表条目（双向） & R153 + lookup\_correct &
witness\_iff\_table\_entry \\
★对偶 & 模糊存在性 ∧ 结构判等 & R147/R152/R153 &
np\_fuzzy\_structure\_duality \\
全景 & N=外推 ∧ 验证=判等 ∧ P=锁定链 & R150/R152/R153/R050 &
p\_np\_pat\_perspective \\
\end{longtable}
}

\textbf{4
个新定理（R157），只锚筑基篇定理（R050/R147/R150/R152/R153/RulerLookup/RulerRevLookup），不用外部引理，\texttt{lake\ build}
通过，0 sorry。}

习题的回答（一句话）：\textbf{NP 的 pat 语义 = 用模糊换速度（N =
相位锁定外推给出存在性），用结构找回精确（验证 = 查表判等还原真值）；P =
锁定方向唯一链，无模糊；有限域上结构已完备，模糊不必要，P = NP
平凡。}这是结构论述，不是 P≠NP
判定------后者需要计算模型的规模增长下界，不在筑基篇能力内。

\begin{center}\rule{0.5\linewidth}{0.5pt}\end{center}

\subsection{六、相关对照}\label{ux516dux76f8ux5173ux5bf9ux7167}

\begin{quote}
完整书目（全字段 + 验证记录 + 健在/已故标注，85+
条，按学科分类）见\texttt{../shared/pat\_pnp\_shared\_references\_ZH.md}（共享引用清单）。正文标注
{[}C\#{]}/{[}M\#{]}/{[}P\#{]}/{[}L\#{]}/{[}Q\#{]}/{[}S\#{]}
对应清单条目。
\end{quote}

{\def\LTcaptype{none} % do not increment counter
\begin{longtable}[]{@{}
  >{\raggedright\arraybackslash}p{(\linewidth - 4\tabcolsep) * \real{0.3333}}
  >{\raggedright\arraybackslash}p{(\linewidth - 4\tabcolsep) * \real{0.3333}}
  >{\raggedright\arraybackslash}p{(\linewidth - 4\tabcolsep) * \real{0.3333}}@{}}
\toprule\noalign{}
\begin{minipage}[b]{\linewidth}\raggedright
文献
\end{minipage} & \begin{minipage}[b]{\linewidth}\raggedright
对应内容
\end{minipage} & \begin{minipage}[b]{\linewidth}\raggedright
备注
\end{minipage} \\
\midrule\noalign{}
\endhead
\bottomrule\noalign{}
\endlastfoot
\textbf{Cook, S. A.} 1971. \emph{The Complexity of Theorem-Proving
Procedures}, STOC {[}C1{]} & NP 完备性（R152/R153 对照基线） &
复杂性理论起点 \\
\textbf{Levin, L. A.} 1973. \emph{Universal Search Problems} {[}C2{]} &
NP 完备性独立发现 & 与 Cook 并列 \\
\textbf{Karp, R. M.} 1972. \emph{Reducibility Among Combinatorial
Problems} {[}C3{]} & NP 完备性 21 问题 & 决策/搜索问题三分 \\
\textbf{Hartmanis, J., Stearns, R. E.} 1965. TAMS {[}C4{]} &
复杂性层级奠基（步数度量） & DOI 已验证 \\
\textbf{Aaronson, S.} 2011. arXiv:1108.1791 {[}C12{]} &
复杂性→哲学：为什么''验证便宜/求解贵''值得追问 &
用模糊换速度的哲学对照 \\
\textbf{Wigderson, A.} 2019. \emph{Mathematics and Computation} {[}C9{]}
& 计算 = 数学结构 & 结构找回精确的纲领性对照 \\
\textbf{Sipser / Papadimitriou / Arora--Barak / Goldreich} {[}C5--C8{]}
& P/NP/验证/表语义的标准形式化 & 教材基线 \\
\textbf{Fortnow, L.} 2013. \emph{The Golden Ticket} {[}C10{]} & P vs NP
的存在性/验证通俗论述 & 对照：验证 vs 求解不对称 \\
\textbf{Valiant, L.} 2013. \emph{Probably Approximately Correct}
{[}C11{]} & 学习 = 计算（模糊先验→精确结构） &
与本题''模糊→精确''同型 \\
\textbf{Chomsky / Pinker / Lakoff / Jackendoff / Goldberg} {[}L1--L6{]}
& 语言 = 结构（语法/构式/语义） & 结构找回精确的语言学对照 \\
\textbf{Berry / Hestenes / Aharonov--Bohm} {[}Q1--Q4{]} & 几何相位 =
相位锁定；几何代数 = 方向声明 & R138 相位锁定的物理对照 \\
\textbf{Gödel / Turing / Church / Kleene} {[}M1--M3, M7{]} &
可计算性/不可判定 = 形式结构边界 & 模糊/精确边界的元对照 \\
\textbf{Shannon / Landauer / Bennett} {[}M10--M12{]} & 信息 =
结构；可逆计算 = 信息保持 & R057 一切皆表的信息论对照 \\
\textbf{de Moura / mathlib / Buzzard / Avigad / Hales / Scholze} {[}M21,
M23, M28--M31{]} & Lean 形式化 = 用结构找回精确的当代实践 & 与本题 Lean
验收同工程 \\
\textbf{R150 王氏定理} & N = 相位锁定外推的存在性来源 &
本框架定理，非外部引理 \\
\end{longtable}
}

\begin{quote}
诚实边界延续 R152：完整 P≠NP
判定需要计算模型下界证明，不在本框架能力内------本题交付的是结构侧写（用模糊换速度、用结构找回精确的详细论述），标注
PROVED 的仅为上述 4 个 Lean
验收定理本身。引用清单纪律：每条均已验证（wiki 标题页 / arXiv abs 页 /
DOI 解析），未验证条目不列入。
\end{quote}

\begin{center}\rule{0.5\linewidth}{0.5pt}\end{center}

\emph{筑基篇课后习题 I · 2026-08-13 · Lean 4 / mathlib v4.32.2 · 4
theorems PROVED no sorry · 配套 Lean
形式化：formal/Formal/Toolkit/PatPvsNPExercise.lean（快照见
../lean/PatPvsNPExercise.lean）}

\end{document}
